% 原书分章 PDF 第 9—12 页：第 25.3.2 节，完整定理与引理证明。
\subsection{CV-TMLE 的渐近正态性}

下面的定理在不对 $P_0$ 作光滑性假设的情况下，沿序列 $h_n\to0$ 建立
$\Psi_{h_n}(P_0)$ 的 CV-TMLE
$\psi^*_{n,h_n}=E_{B_n}\Psi_{h_n}(\mathbf P^*_{n,B_n,h_n})$
的渐近正态性。

\begin{bookresult}{定理}{25.1}
作如下假设。
\begin{itemize}
\item 令 $r(n)$ 为趋于零的速率，使对某个
$P_1\in\mathcal M$，所有 $h>0$ 都有 $R_h(P_1,P_0)=0$，且
\begin{equation}
 \|D_h^*(\mathbf P^*_{n,B_n,h})-D_h^*(P_1)\|_{P_0}
 =O_P(h^{-d/2}r(n)).
\tag{25.3}
\end{equation}
\item 令 $h_n\to_p0$，
$r(n)\{1-\log(h_n^{d/2}r(n))\}\to0$，且
\begin{equation}
 \lim_{h\to0}h^dP_0\{D_h^*(P_1)\}^2=\sigma_0^2.
\tag{25.4}
\end{equation}
\item 存在确定性序列 $\tilde h_n$，使 $h_n/\tilde h_n\to_p1$。
\item 考虑函数类
\[
 \mathcal F_n=
 \{D(K_{1,h,x_0},P_{n,B_n,\epsilon,h})
  -D(K_{1,h,x_0},P_1):h,\epsilon\},
\]
并令 $F_n$ 为其包络函数。假定对某个整数 $p>0$，
\[
 \sup_\Lambda N(\delta|F_n|,\mathcal F_n,L^2(\Lambda)
 =O(\delta^{-p}).
\]
\item
\begin{equation}
 \sqrt{nh_n^d}\,E_{B_n}\int_y
 K(y)R_{x_0+h_ny}(\mathbf P^*_{n,B_n,h_n},P_0)\,dy=o_P(1).
\tag{25.5}
\end{equation}
\end{itemize}
则
\[
\begin{aligned}
 &\sqrt{nh_n^d}\bigl(
 E_{B_n}\Psi_{h_n}(\mathbf P^*_{n,B_n,h_n})-\Psi_{h_n}(P_0)\bigr)\\
 &\quad=\sqrt{nh_n^d}(P_n-P_0)D_{h_n}^*(P_1)+o_P(1)\\
 &\qquad+\sqrt{nh_n^d}\,E_{B_n}\int_y
 K(y)R_{x_0+h_ny}(\mathbf P^*_{n,B_n,h_n},P_0)\,dy,
\end{aligned}
\]
且
\[
 \sqrt{nh_n^d}\bigl(
 E_{B_n}\Psi_{h_n}(\mathbf P^*_{n,B_n,h_n})-\Psi_{h_n}(P_0)\bigr)
 \Rightarrow_dN(0,\sigma_0^2)).
\]
\end{bookresult}

回顾 $D_h^*(P)=D(K_{h,x_0},P)$，其中 $K\mapsto D(K,P)$ 关于 $K$ 线性。
例如，$O=(X,O_1)$，且
\[
 D_h^*(P)=K_{h,x_0}D_1(P)
 +\int K_{h,x_0}(X)D_2(P)(x,O_1)\,dx,
\]
其中
\[
 \int K_{h,x_0}(x)D_2(P)(x,O_1)\,dx
 =\int K(y)D_2(P)(x_0+hy,O_1)\,dy.
\]
此时，式 (25.3) 左端通常由
$\|K_{h,x_0}D_1(\mathbf P^*_{n,B_n,h})-K_{h,x_0}D_1(P_1)\|_{P_0}$
控制，可选择 $r(n)$ 使
\[
\begin{aligned}
 \biggl(\int_y K^2(y)E_0\bigl(
 &\{D_1(\mathbf P^*_{n,B_n,h_n})-D_1(P_1)\}^2
 \mid X=x_0+h_ny\bigr)\\
 &\quad p_{0,X}(x_0+h_ny)\,dy\biggr)^{0.5}
 =O_P(r(n)).
\end{aligned}
\]

满足 $R_{20}(P_1,P_0)=0$ 的最简单情形是 $P_1=P_0$。
在许多问题中，有双重稳健结构
\[
\begin{aligned}
 R_x(P,P_0)=\int_y
 &(H_1(P)-H_1(P_0))(x,y)\\
 &\quad(H_2(P)-H_2(P_0))(x,y)H_3(P,P_0)(x,y)\,d\mu(y).
\end{aligned}
\]
此时，只要 $H_1(P_1)=H_1(P_0)$ 或 $H_2(P_1)=H_2(P_0)$ 之一成立，
$R_x(P,P_1)$ 就等于零。例如，考虑 $H_2(P_1)=H_2(P_0)$ 的情形。有如下分解：
\[
\begin{aligned}
 R_x(\mathbf P_n,P_0)
 &=\int_y(H_1(\mathbf P_n)-H_1(P_1))
   (H_2(\mathbf P_n)-H_2(P_0))(x,y)\\
 &\qquad\qquad H_3(\mathbf P_n,P_0)(x,y)\,d\mu(y)\\
 &\quad+\int_y(H_1(P_1)-H_1(P_0))
   (H_2(\mathbf P_n)-H_2(P_0))(x,y)\\
 &\qquad\qquad H_3(\mathbf P_n,P_0)(x,y)\,d\mu(y).
\end{aligned}
\]
因此，若右端两项都为 $o_P((nh^d)^{-0.5})$，就满足条件 (25.9)。
这将要求 $H_2(\mathbf P_n)$ 以快于 $(nh^d)^{-0.5}$ 的速率收敛到 $H_2(P_0)$。

\begin{proof}[定理 25.1 的证明]
我们有
\[
\begin{aligned}
 &\sqrt{nh_n^d}\bigl(
 E_{B_n}\Psi_{h_n}(\mathbf P^*_{n,B_n,h_n})-\Psi_{h_n}(P_0)\bigr)\\
 &\quad=\sqrt{nh^d}\,E_{B_n}(P^1_{n,B_n}-P_0)
 D_{h_n}^*(\mathbf P^*_{n,B_n,h_n})\\
 &\qquad+\sqrt{nh_n^d}\,E_{B_n}
 R_{h_n}(\mathbf P^*_{n,B_n,h_n},P_0).
\end{aligned}
\]
考虑主导经验过程项。令
$f_{n,B_n}\equiv
D_{h_n}^*(\mathbf P^*_{n,B_n,h_n})-D_{h_n}^*(P_1))$。
令 $K_{1,h,x_0}(x)=K((x-x_0)/h)$，从而
$K_{h,x_0}=h^{-d}K_{1,h,x_0}$。
注意，$f_{n,B_n}=h_n^{-d}g_{n,B_n}$，其中
$g_{n,B_n}=D(K_{1,h_n,x_0},P^*_{n,B_n,h_n})
-D(K_{1,h_n,x_0},P_1)$。
令 $p=\sum_{i=1}^{n}B_n(i)/n$。有
\[
\begin{aligned}
 &\sqrt{nh_n^d}\,E_{B_n}(P^1_{n,B_n}-P_0)
 \bigl(D_{h_n}^*(\mathbf P^*_{n,B_n,h_n}-D_{h_n}^*(P_1)\bigr)\\
 &\quad=\sqrt{nh_n^d}\,E_{B_n}(P^1_{n,B_n}-P_0)f_{n,B_n}\\
 &\quad=\sqrt{nh_n^d}\,h_n^{-d}E_{B_n}
 (P^1_{n,B_n}-P_0)g_{n,B_n}\\
 &\quad=p^{-0.5}h_n^{-0.5d}\sqrt{np}\,E_{B_n}
 (P^1_{n,B_n}-P_0)g_{n,B_n}.
\end{aligned}
\]
对于给定的 $B_n$，以 $P^0_{n,B_n}$ 为条件，将下面的引理 25.2 应用于
$\sqrt{np}(P^1_{n,B_n}-P_0)g_{n,B_n}$。
其中，$g_{n,B_n}$ 仅通过 $\theta_n=(\epsilon_n,h_n)$ 具有随机性。
为简化记号，将该函数 $g_{n,B_n}$ 记为 $g_{n,\theta_n}$。
由于 $g_{n,\theta_n}\in\mathcal F_n=\{g_{n,\theta}:\theta\}$，
$\mathcal F_n$ 的确如引理 25.2 中那样是有限维的。
此外，根据 $D_{h_n}^*(P)$ 关于 $o$ 和 $P\in\mathcal M$ 一致地被
$Ch_n^{-d}$ 控制的假设，$\mathcal F_n$ 中任意函数都被统一常数 $C$ 控制。
因此，对于某个 $C<\infty$，$\mathcal F_n$ 的包络函数 $F_n$ 满足
$\|F_n\|_{P_0}<C<\infty$。
由假设，$\|g_{n,\theta_n}\|_{P_0}=O_P(h_n^{d/2}r(n))$。
应用引理 25.2 得
\[
 \sqrt{np}(P^1_{n,B_n}-P_0)g_{n,\theta_n}
 =O_P\bigl(h_n^{d/2}r(n)\{1-\log(h_n^{d/2}r(n))\}\bigr).
\]
由于 $B_n$ 只有有限多个实现，这证明了
\[
 \sqrt{np}\,E_{B_n}(P^1_{n,B_n}-P_0)g_{n,B_n}
 =O_P\bigl(h_n^{d/2}r(n)\{1-\log(h_n^{d/2}r(n))\}\bigr).
\]
由此可见，
\[
\begin{aligned}
 &p^{-0.5}h_n^{-0.5d}\sqrt{np}\,E_{B_n}
 (P^1_{n,B_n}-P_0)g_{n,B_n}\\
 &\qquad=O_P\bigl(r(n)\{1-\log(h_n^{d/2}r(n))\}\bigr).
\end{aligned}
\]
由假设，右端为 $o_P(1)$。因此，我们已证明
\[
\begin{aligned}
 &\sqrt{nh_n^d}\,E_{B_n}(P^1_{n,B_n}-P_0)
 \bigl(D_{h_n}^*(\mathbf P^*_{n,B_n,h_n})-D_{h_n}^*(P_1)\bigr)\\
 &\qquad=o_P\bigl(r(n)\{1-\log(h_n^{d/2}r(n))\}\bigr)=o_P(1).
\end{aligned}
\]
现在得到
\[
\begin{aligned}
 &\sqrt{nh_n^d}\bigl(
 E_{B_n}\Psi_{h_n}(\mathbf P^*_{n,B_n,h_n})-\Psi_{h_n}(P_0)\bigr)\\
 &\quad=\sqrt{nh_n^d}\,E_{B_n}(P^1_{n,B_n}-P_0)D_{h_n}^*(P_1)+o_P(1)\\
 &\qquad+\sqrt{nh_n^d}\,E_{B_n}\int_x K_{h_n,x_0}(x)
 R_x(\mathbf P^*_{n,B_n,h_n},P_0)\,dx.
\end{aligned}
\]
利用假设 (25.4)，由三角阵列的 Lindeberg 定理可知，对于确定性序列 $\tilde h_n$，
\[
 \sqrt{n\tilde h_n^d}\,E_{B_n}(P^1_{n,B_n}-P_0)
 D_{\tilde h_n}^*(P_1)\Rightarrow_dN(0,\sigma_0^2).
\]
由假设，存在确定性序列 $\tilde h_n$，使 $h_n/\tilde h_n\to_p1$。有
\[
\begin{aligned}
 &\sqrt{nh_n^d}\,E_{B_n}(P^1_{n,B_n}-P_0)D_{h_n}^*(P_1)\\
 &\quad-\sqrt{n\tilde h_n^d}\,E_{B_n}(P^1_{n,B_n}-P_0)
 D_{\tilde h_n}^*(P_1)\\
 &=n^{0.5}\{\sqrt{h_n^d}-\sqrt{\tilde h_n^d}\}
 E_{B_n}(P^1_{n,B_n}-P_0)D_{\tilde h_n}^*(P_1)\\
 &\quad+\sqrt{nh_n^d}\,E_{B_n}(P^1_{n,B_n}-P_0)
 \{D_{h_n}^*(P_1)-D_{\tilde h_n}^*(P_1)\}.
\end{aligned}
\]
第一项可写为 $\sqrt{h_n^d/\tilde h_n^d}-1$ 乘以
$(n\tilde h_n^d)^{0.5}E_{B_n}(P^1_{n,B_n}-P_0)D_{\tilde h_n}^*(P_1)$。
由于该标量趋于零，且经验过程项收敛到 $N(0,\sigma_0^2)$，第一项为 $o_P(1)$。
第二项可按上述经验过程项的方法分析，因此也为 $o_P(1)$。
现在考虑二阶余项
$\sqrt{nh_n^d}\,E_{B_n}R_{h_n}(\mathbf P^*_{n,B_n,h_n},P_0)$。
它等于
\[
 \sqrt{nh_n^d}\,E_{B_n}\int_y K(y)
 R_{x_0+hy}(\mathbf P^*_{n,B_n,h},P_0)\,dy.
\]
由假设 (25.5)，后者为 $o_P(1)$。定理得证。
\end{proof}

证明中使用了以下引理。

\begin{bookresult}{引理}{25.2}
令 $G_n(f)=\sqrt n(P_n-P_0)f$ 为 $f(O_i)$ 的标准化经验均值。
令 $f_{n,\epsilon_n}\in\mathcal F_n=\{f_{n,\epsilon}:\epsilon\}$，
其中 $\epsilon$ 在 $\mathbb R^p$ 的一个有界集合中变化，$f_{n,\epsilon}$ 为非随机函数，即不基于数据 $O_1,\ldots,O_n$。
假定 $\|f_{n,\epsilon_n}\|_{P_0}=o_P(r_n)$，
其中速率 $r_n$ 满足 $r_n\log r_n^{-1}\to0$ 和 $r_nn^{0.5}\to\infty$。
假定 $\mathcal F_n$ 的包络函数 $F_n$ 满足 $\|F_n\|_{P_0}\leq M_n$。
假定对某个整数 $p>0$，
\[
 \sup_\Lambda N(\delta|F_n|,\mathcal F_n,L^2(\Lambda)=O(\delta^{-p}).
\]
则
\[
 E_0|G_n(f_{n,\epsilon_n})|
 =O\bigl(M_nr_n(1+\log r_n^{-1})\bigr).
\]
因此，若 $r_n=o(1)$ 且 $M_n<M<\infty$，则
$G_n(f_{n,\epsilon_n})=o_P(1)$。
\end{bookresult}

\begin{proof}
为简化记号，将 $f_{n,\epsilon_n}$ 记为 $f_n$。
应用 van der Vaart 和 Wellner（2011）的定理，得到
\begin{equation}
 E_0|G_n(f_n)|
 \leq J(\delta_n,\mathcal F_n)
 \left(1+\frac{J(\delta_n,\mathcal F_n)}
 {\delta_n^2n^{0.5}\|F_n\|_{P_0}}\right)\|F_n\|_{P_0},
\tag{25.6}
\end{equation}
其中可选择 $\delta_n=r_n$。
利用一致覆盖数的上界 $\epsilon^{-p}$，可得
$J(\delta_n,\mathcal F_n)
=-p^{0.5}\int_0^{\delta_n}(\log\epsilon)^{0.5}\,d\epsilon$。
我们可以保守地用 $\log\epsilon$ 控制 $(\log\epsilon)^{0.5}$，并利用
\[
 \int_0^{\delta_n}\log\epsilon\,d\epsilon
 =\delta_n-\delta\log\delta_n
 =\delta(1+\log(\delta_n^{-1})).
\]
这表明 $J(\delta,\mathcal F_n)\leq\delta_n(1+\log\delta_n^{-1})$。
若 $J(\delta_n,\mathcal F_n)=O(\delta_n^2n^{0.5})$，
则式 (25.6) 的主导项为
$J(\delta_n,\mathcal F_n)\|F_n\|_{P_0}$。
利用上述 $J(\delta_n,\mathcal F_n)$ 的界可知，只要
$\delta_n(1+\log\delta_n^{-1})=O(\delta_n^2n^{0.5})$，
或等价地，$\delta_n(1+\log\delta_n^{-1})=O(\delta_nn^{0.5})$，
这一点就成立。由假设，
$\delta_nn^{0.5}\to\infty$ 且 $\delta_n\log\delta_n^{-1}\to0$，
因此它总是成立。由此得到
\[
 E_0|G_n(f_n)|=O(M_nr_n(1+\log r_n^{-1})),
\]
即所述的界。
\end{proof}
